\documentclass{beamer}

\usepackage{amsmath,amssymb,amsthm}
\usepackage[all]{xy}



\numberwithin{equation}{subsection}

% Basic macros
\def\lim{\mathop{\mathrm{lim}}\nolimits}
\def\colim{\mathop{\mathrm{colim}}\nolimits}
\def\Spec{\mathop{\mathrm{Spec}}}
\def\Hom{\mathop{\mathrm{Hom}}\nolimits}
\def\Ext{\mathop{\mathrm{Ext}}\nolimits}
\def\Ker{\mathop{\mathrm{Ker}}}
\def\Im{\mathop{\mathrm{Im}}}
\def\Coker{\mathop{\mathrm{Coker}}}
 \def\Coim{\mathop{\mathrm{Coim}}}

% Slide numbering
\addtobeamertemplate{navigation symbols}{}{%
    \usebeamerfont{footline}%
    \usebeamercolor[fg]{footline}%
    \hspace{1em}%
    \insertframenumber/\inserttotalframenumber
}



\title{Nakajima's representation of the Heisenberg algebra}
\author{Daniel González Casanova Azuela}
\institute{PUC-Rio}
\date{\today}

\begin{document}

\begin{frame}%Title page
\titlepage
\end{frame}

\section*{Table of contents}
\label{section-toc}

\begin{frame}%Plan of the talk
\frametitle{Plan of the talk}
\tableofcontents
\end{frame}

\section{Heisenberg Algebra}

\begin{frame}%Heisenberg algebra
\frametitle{Heisenberg algebra}
The Heisenberg algebra is
$
\hat{\mathfrak{a}}=
\bigoplus_{n \in \mathbb{Z}} \mathbb{C}h_n \oplus\mathbb{C}K
$
with bracket given by 
$$
[h_m,h_n]=m\delta_{m,-n}K,\qquad 
[K,\hat{\mathfrak{a}}]=0,
$$
and highest weight representation 
$$
H=\mathbb{C}[x_1,x_2,\ldots]$$
with the action
$$
h_n\mapsto \begin{cases}
n \frac{\partial }{\partial x_n}\qquad &\text{if }n>0\\
x_{-n}\qquad &\text{if }n<0\\
0\qquad &\text{if }n=0
\end{cases},
\qquad \qquad  K\mapsto \text{Id}.
$$
\end{frame}

\section{Hilbert Scheme of Points}

\begin{frame}%Hilbert scheme of points
\frametitle{Hilbert scheme of points}
Let $X$ be a projective surface over $\mathbb{C}$.
The essential property of the 
{\it Hilbert scheme
of $n$ points $X^{[n]}$}
is that a point (over $\mathbb{C}$) of $X^{[n]}$
is a 0-dimensional closed subscheme of $X$
of length $n$.

\medskip\noindent
\pause
The {\it length} of any such subscheme $Z$, defined by
 $\dim\Gamma(\mathcal{O}_Z)$,
is given by
$\sum_{x \in \text{supp} Z}\dim\Gamma(\mathcal{O}_{Z_x})=n$.

\medskip\noindent
\pause
This {\bf does not} mean that the points in $X^{[n]}$ 
are $n$-tuples of points of $X$ counted with multiplicities.
\end{frame}

\begin{frame}%Fogarty theorem
\frametitle{Hilbert scheme of points}
\begin{theorem}[Fogarty]
\label{theorem-fogarty}
If $X$ is a nonsingular projective surface,
\begin{enumerate}
\item $X^{[n]}$ is nonsingular and has dimension $2n$, and
\item the {\it Hilbert-Chow morphism}
\begin{align*}
\pi : X^{[n]} &\longrightarrow S^nX \\
 Z&\longmapsto \sum_{x \in \text{supp}Z}\dim\Gamma(\mathcal{O}_{Z,x})[x]
\end{align*}
is a resolution of singularities.
\end{enumerate}
\end{theorem}
\end{frame}

\section{Borel-Moore Homology}

\begin{frame}%Borel-Moore homology
\frametitle{Borel-Moore homology}
{\it Borel-Moore homology $H_\bullet^{BM}(X)$}
is the homology of the chain complex
of {\bf locally finite} linear combinations
of singular simplices.

\pause
\begin{enumerate}
\item 
\label{item-pushforward}
(Pushforward.)
If $f:X \to Y$ is proper
then there is a map $f_* :H_i^{BM}(X)\to H_i^{BM}(Y)$.

\pause
\item 
\label{item-pullback}
(Pullback.) If $f:X \to Y$ 
is a locally trivial fibration with
fiber a manifold of dimension $m$ (over $\mathbb{R}$),
then there is a map
$f^* :H_i^{BM}(Y)\to H_{i+m}^{BM}(X)$.

\pause
\item  
\label{item-intersection}
(Intersection pairing.) If $X$ and $Y$ 
are both ``nicely embedded'' in the same
manifold $M$, there is a map
$$
\cap_M:H_i^{BM}(X)\times H_j^{BM}(Y)
\to H_{i+j-m}^{BM}(X\cap Y).
$$
\vspace{-1.5em}\pause
\item 
\label{item-fundamental-class}
(Fundamental class.)
If $X$ is a variety of dimension $d$ over $\mathbb{C}$ then
there is {\it fundamental class} $[X] \in H^{BM}_{2d}(X)$.
\end{enumerate}
\end{frame}

\section{Correspondences}

\begin{frame}%Correspondences
\frametitle{Correspondences}
Let $X_1,X_2$ be smooth irreducible varieties
of dimensions $d_1,d_2$.
Let $Z$ be a $d$-dimensional 
closed subvariety of $X_1\times X_2$ 
such that the projections $p_1$ and $p_2$
are proper.

\medskip\noindent
The {\it correspondence} associated to $Z \subset X_1\times X_2$
is the map
\begin{align*}
c_{[Z]}:H_i^{BM}(X_2)&\longrightarrow H_{i+2d-2d_2}^{BM}(X_1)\\
c&\longmapsto p_{1,*}(p_2^*c\cap [Z])
\end{align*}

\noindent
where $p_1$ and $p_2$ are the projections in the following diagram:
$$
\xymatrix{
&Z\ar[ld]_{p_1}\ar[rd]^{p_2} \subset X_1\times X_2\\
X_1&&X_2.
}
$$
\end{frame}

\begin{frame}%Composition of correspondences
\frametitle{Composition of correspondences}
Let $Z_{12} \subset X_1 \times X_2$ and $Z_{23} \subset X_2 \times X_3$
such that the projections to $X_2$ and $X_3$ are proper.
The composition of correspondences is
$$
[Z_{12}]\circ[Z_{23}]:=p_{13,*}(p_{12}^* [Z_{12}]\cap p_{23}^* [Z_{23}]),
$$
$$
\xymatrix{
&  X_1 \times X_2 \times X_3\ar[dl]_{p_{12}}\ar[dr]^{p_{23}}\ar[dd]_{p_{13}}\\
 X_1\times X_2& & X_2\times X_3\\
& X_1\times X_3.
}
$$
In fact,
the projection $p_{13}(p_{12}^{-1}Z_{12} \cap p_{23}^{-1} Z_{23})\to X_3$ 
is proper.
\end{frame}


\section{Main Construction}

\begin{frame}%Main construction
\frametitle{Main construction}
Let $i>0$. Define $P[i] \subseteq
\coprod_nX^{[n-i]}\times X^{[n]}\times X$
by
$$
P[i]:=\coprod_n 
\left\{(\mathcal{J}_1\supseteq\mathcal{J}_2,x):
\text{supp}(\mathcal{J}_1/\mathcal{J}_2)=\{x\}\right\}.
$$

\pause
Let
$\alpha \in H_*^{BM}(X)$, $\beta \in H_*(X)$
and $i>0$. Define
\begin{align*}
P_\alpha[i]&=\varpi_*(\Pi_1^*\alpha \cap [P[i]])\\
P_\beta[i]&=\varpi_*(\Pi_1^*\beta \cap [P[-i]]).
\end{align*}
$$
\xymatrix{
&  \coprod_n X^{[n-i]}\times X^{[n]}\times X
\ar[dl]^{\Pi_1}\ar[dr]_{\varpi_i}\\
X&&\coprod_n X^{[n-i]}\times X^{[n]}.
}
$$
\end{frame}

\begin{frame}%Main theorem
\frametitle{Main theorem}
\begin{theorem}[Nakajima, Grojnowski]
\label{theorem-main}
If $(\text{deg} \alpha)(\text{deg}\beta)\equiv 0 \text{ mod } 2$,
$$
[P_\alpha[i],P_\beta[j]]=(-1)^{i-1}i \delta_{i,-j}
\left<\alpha,\beta\right>\text{Id},
$$
where $\left<\alpha,\beta\right>=p_*(\alpha \cap \beta)$
for $p:X \to \text{pt}$ the unique morphism
from $X$ to the point pt.
\end{theorem}
\end{frame}

\section{Idea of Proof}

\begin{frame}%Easy case
\frametitle{Easy case}

Let $x \in X$, $\alpha=[X]$, $\beta=[x]$.
Take a set of distinct $n$ points
$x_1,\ldots,x_n \in X\setminus\{x\}$ 
and identify them with the
ideal $\mathcal{J}$. Then
$$
[P_\alpha[1],P_\beta[-1]][\mathcal{J}]=[\mathcal{J}].
$$

\pause\vspace{-1em}
$$
\xymatrix{
&X^{[n+1]}\times X^{[n]}\ar[ld]_{p_1}\ar[rd]^{p_2}\\
X^{[n+1]}&&X^{[n]}.
}
$$

\begin{itemize}
\item $P_{[x]}[-1][\mathcal{J}]=[\{x\}\cup \mathcal{J}]$,
since we are looking at
$$
\{(\mathcal{J'}\supseteq\mathcal{J},x):
\text{supp}(\mathcal{J}'/\mathcal{J})=\{x\}\}
\subset X^{[n]} \times X^{[n+1]} \times X.
$$

\item $P_{[X]}[1]P_{[x]}[-1][\mathcal{J}]$, since
$$
\{(\mathcal{J}''\subseteq \mathcal{J}',x)
:\text{supp}(\mathcal{J}'/\mathcal{J}'')=\{y\},
\text{ for some }y\in X\}
$$
\end{itemize}
\end{frame}

\begin{frame}%Easy case
\frametitle{Easy case}

\begin{center}
\includegraphics[width=0.9\textwidth]{fig1.png}
\end{center}

\end{frame}

\begin{frame}%Proof for i,j>0
\frametitle{Proof for $i,j>0$}

Let's explain the case $i,j>0$.
By definition,
$$
P_\alpha[i]P_\beta[j]
=p_{13,*}(p_{12}^*[P[i]]\cap\Pi_1^*\alpha\cap p_{23}^*[P[j]]\cap\Pi_2^*\beta).
$$
$$
\xymatrixcolsep{.2em}
\xymatrix{
&  X^{[n-i-j]} \times X^{[n-j]} \times X^{[n]}
\ar[dl]_{p_{12}}\ar[dr]^{p_{23}}\ar[dd]_{p_{13}}\\
P[j]\subset X^{[n-j]}\times X^{[n]}& & X^{[n-i-j]}\times X^{[n-j]}\supset P[i]\\
& X^{[n-i-j]}\times X^{[n]}
}
$$
\end{frame}

\begin{frame}%Proof for i,j>0
\frametitle{Proof for $i,j>0$}
\begin{align*}
P[i,j]&:=p_{12}^{-1}(P[i])\cap p_{23}^{-1}(P[j])
\cap(X^{[n-i-j]}\times X^{[n-j]}\times X^{[n]})\\
&=\left\{
(\mathcal{J}_1\supseteq\mathcal{J}_2\supseteq\mathcal{J}_3):
\substack{
\text{supp}(\mathcal{J}_1/\mathcal{J}_2)=\{x\}
\text{ for some }x\in X\\
\text{supp}(\mathcal{J}_2/\mathcal{J}_3)=\{y\}
\text{ for some }y \in X }\right\}
\end{align*}

\pause
$P[i,j]$ is the disjoint union of
of $\{x \neq y\}$ and $\{x=y\}$.

\medskip\noindent
Let $\overline{P}[i,j]:=\overline{P[i,j]\cap\{x\neq y\}}$.
Then
$$
p_{12}^* [P[i]]\cap p_{23}^* [P[j]]=[\overline{P}[i,j]]+\iota_*R
$$
where $R \in H_*^{BM}(P[i,j]\cap \{x=y\})$ corresponds
to the $\{x=y\}$ part of $P[i,j]$ 
and $\iota$ is the inclusion.
\end{frame}

\begin{frame}%Lemma
\frametitle{Proof for $i,j>0$}
\begin{lemma}
\label{lemma-pij}
\begin{enumerate}
\item There exists an isomorphism
$$
P[i,j]\cap\{x \neq y\}\xrightarrow{\cong}P[j,i]\cap \{x \neq y\},
$$
\noindent
which exchanges $\Pi_1$ and $\Pi_2$.

\item $
p_{13,*}(\Pi_1^*\alpha\cap\Pi_2^*\beta\cap\iota_*R)=0,\qquad 
p_{13,*}(\Pi_1\beta\cap\Pi_2^*\alpha\cap\iota_*R')=0.
$
\end{enumerate}
\end{lemma}

\medskip\noindent\pause
\noindent
This lemma implies the result since
\begin{align*}
P_\alpha[i]P_\beta[j]
&=p_{13,*}(\overline{P}[i,j]\cap\Pi_1^*\alpha\cap\Pi_2^*\beta)\\
&=(-1)^{\text{deg}\alpha\text{deg}\beta}
p_{13,*}(\overline{P}[j,i]\cap\Pi_1^*\beta\cap\Pi_2^*\alpha)\\
&=P_\beta[j]P_\alpha[i].
\end{align*}
\end{frame}

\begin{frame}%Proof of 1.
\frametitle{Proof of 1.}
Define an isomorphism by
$$
\mathcal{J}_1/\mathcal{J}_2'=\mathcal{J}_2/\mathcal{J}_3,\qquad
\mathcal{J}_2'/\mathcal{J}_3=\mathcal{J}_1/\mathcal{J}_2.
$$
\begin{center}
\includegraphics[width=0.8\textwidth]{fig2.png}
\end{center}

\end{frame}

\begin{frame}%Proof of 2.
\frametitle{Proof of 2.}
Consider the commutative diagram
$$
\xymatrix{
P[i,j]\cap \{x=y\}\ar[r]^-{\Pi'}\ar[d]_{\iota}&X\ar[d]^{\Delta}\\
P[i,j]\ar[r]_{\Pi_1 \times \Pi_2}&X \times X.
}
$$
where $\Delta$ is the diagonal embedding
and $\Pi'$ is any of the projections $\Pi_1$ or $\Pi_2$.
Then
\begin{align*}
\Pi_1^*\alpha \cap \Pi_2^*\beta \cap\iota_*R
&=\iota_*(\iota^*(\Pi_1^*\alpha\cap\Pi_2^*\beta)\cap R)\\
&=\iota_*(\iota^*(\Pi_1\times\Pi_2)^*(\alpha\times\beta)\cap R)\\
&=\iota_*({\Pi'}^*\Delta^*(\alpha\times\beta)\cap R)\\
&=\iota_*({\Pi'}^*(\alpha\cap\beta)\cap R).
\end{align*}
\end{frame}

\begin{frame}%Proof of 2.
\frametitle{Proof of 2.}
Consider another commutative diagram:
$$
\xymatrix{
P[i,j]\cap\{x=y\}
\ar[r]^-{\Pi'}\ar[d]_{p_{13 \circ \iota}}&X\ar@{=}[d]\\
P[i+j]=\{(\mathcal{J}_1\supseteq\mathcal{J}_3):
\text{supp}(\mathcal{J}_1/\mathcal{J}_3)=\{x\}\}\ar[r]_-{\Pi''}&X,
}
$$
where $\Pi''$ is the usual projection to $X$ from the
set $P[i+j]$. Then
$$
p_{13,*}\iota_*({\Pi'}^* (\alpha\cap \beta)\cap R)
={\Pi''}^* (\alpha \cap \beta)\cap p_{13,*}\iota_*(R).
$$
Since $P[i+j]$ has less than the ``expected dimension'',
$p_{13,*}\iota_*(R)=0$.
\end{frame}

\begin{frame}%Bibliography
\frametitle{Bibliography}
\begin{thebibliography}{9}
\bibitem{nakajima-lectures} Nakajima, H. Lectures on Hilbert schemes of points on surfaces. In: Algebraic Geometry, Santa Cruz 1995, Proc. Sympos. Pure Math., vol. 62, AMS, 1997, pp. 33--104.
\bibitem{nakajima-paper} Nakajima, H. Heisenberg algebra and Hilbert schemes of points on projective surfaces. Ann. of Math. (2) 145 (1997), no. 2, 379--388.
\bibitem{fogarty} Fogarty, J. Algebraic families on an algebraic surface. Amer. J. Math. 90 (1968), 511--521.
\bibitem{fulton} Fulton, W. Young Tableaux. Cambridge University Press, 1997.
\bibitem{videos} Online videos: \url{https://www.youtube.com/playlist?list=PLtmvIY4GrVv_tNl9B41l6k8Xllk_RGwG0}
\end{thebibliography}
\end{frame}

\end{document}
